\section{Conclusions and policy implications}
\label{sec:conclusions}

We set out to establish which uncertainty governs the trajectory of the Peruvian
power system, rather than to forecast that trajectory. A system dynamics model with
explicit generation and transmission pipelines, coupled to a congestion-aware
dispatch and subjected to a 129-run designed experiment over 2025--2050, gives an
unambiguous answer.

\paragraph{Demand growth dominates, by an order of magnitude.} Over the sampled
ranges, demand growth reaches Spearman rank correlations of $0.985$ with required
transmission, $0.973$ with installed capacity and $-0.752$ with the average price,
while no other factor exceeds $0.27$ in absolute value. A $\pm$20\,\% band on the
2050 demand level moves required transmission by $-18.9$\,\% to $+23.8$\,\% and
installed capacity by $-24.0$\,\% to $+29.3$\,\%. The one-at-a-time and the
all-factors-varying evidence agree on this ranking.

\paragraph{Fuel-price risk is weaker than the literature on Peru implies, and for a
structural reason.} Fuel prices are the weakest driver on three of four indicators,
moving the 2050 average price by \SI{0.005}{\USD\per\mega\watt\hour}. The model's own
expansion path removes the exposure: thermal generation falls to
\SI{7.5}{\tera\watt\hour\per\year} of \SI{148.8}{\tera\watt\hour\per\year} by 2050.
Fuel-price risk in Peru is therefore better understood as a near-term exposure
contingent on transmission delivery than as an independent long-run risk.

\paragraph{The least-cost trajectory is the most network-intensive.} Faster demand
growth lowers the long-run average price, because it sustains the price signal long
enough to pull forward low-cost wind and solar. The best case on delivered price
requires \SI{58562}{\kilo\meter} of transmission, 26.4\,\% more than the central case
and 50\,\% more than the worst, with 10.3\,\% more inter-zonal transfer capability.
For a system whose grid is already the binding constraint, the cheap path is the one
that demands most of the planner.

\paragraph{Lead times change when, not how much.} Transmission lead times rank last
or second-to-last on terminal indicators, but the decomposition by transmission
category shows why: the model builds much the same network either way, and the
deferral surfaces as inter-zonal congestion rather than as a different system size.
Terminal stocks are the wrong indicator for that effect, which is a limitation of
our reporting and a clear direction for the next study.

\paragraph{The expansion is renewable without being told to be.} The central case
reaches 84.9\,\% non-thermal capacity by 2050 under a profitability rule alone, with
no target, mandate or carbon price. That a behavioural model and the
optimisation-based Peruvian literature \citep{fiestas2024,delacruz2024} agree on the
direction of the transition, having disagreed on almost everything methodological, is
itself worth recording.

\subsection*{Policy implications}

The practical reading is that Peruvian planning attention is not proportionate to
the sensitivities. Effort devoted to fuel-price exposure buys little; effort devoted
to (i) long-horizon demand forecasting, (ii) the institutional capacity to deliver
transmission at the rate demand requires, and (iii) moving the network trigger from
confirmed to anticipated generation projects \citep{sauma2006,herding2024}, buys a
great deal. Restoring a predictable procurement signal matters for the same reason:
the investment response in the model is to a sustained price signal, and the
suspension of the auction mechanism removes it.

\subsection*{Limitations}

The ranking is conditional on the ranges in Table~\ref{tab:factors}; a wider delay
distribution, such as the multi-year slippage documented for individual projects
elsewhere, could change it. Storage and grid-enhancing technologies are disabled in
every run reported here, so the substitution between flexibility and network is
absent and the transmission requirement of the high-demand case is, if anything,
overstated. Retirements are exogenous, the dispatch uses one representative day per
month, and the distributions are uniform ranges of admissibility rather than
estimated densities. Each of these is documented with its consequence in the
supplementary material.

\subsection*{Further work}

Three extensions follow directly. First, enabling storage and grid-enhancing
technologies and re-running the same design would quantify how much network the
flexibility options actually displace --- our results identify the Z2--Z3 corridor
under high demand as where the gain would be largest. Second, replacing terminal
stocks with congestion-hours, curtailment and unserved energy as reported indicators
would capture the timing effect of lead times that our indicators miss. Third, the
same design applied to Brazil and Chile with the same model structure would separate
what is Peruvian from what is a property of longitudinal systems generally; the
model and driver used here are already structured for that comparison.
